% 2024年北京高考数学第17题解析图1：空间直角坐标系
\begin{tikzpicture}[>=stealth,line width=0.6pt,font=\normalsize]
  \coordinate (A) at (0,0);
  \coordinate (B) at (0.6,-0.6);
  \coordinate (C) at (2.4,-0.6);
  \coordinate (D) at (5.2,0);
  \coordinate (P) at (1.7,3.5);
  \coordinate (E) at (1.7,0);
  \coordinate (F) at (1.7,1.8);
  \coordinate (S) at (3.5,1.75);

  % 坐标轴 (以E为原点)
  \draw[->] (E) -- ($ (E)!1.3!(C) $) node[below right] {$x$};
  \draw[->] (E) -- ($ (E)!1.2!(D) $) node[below] {$y$};
  \draw[->] (E) -- (1.7,4.2) node[left] {$z$};

  % 实线可见边
  \draw (A) -- (B) -- (C) -- (D);
  \draw (P) -- (A);
  \draw (P) -- (B);
  \draw (P) -- (C);
  \draw (P) -- (D);
  \draw (C) -- (S);

  % 虚线不可见底边与内部线
  \draw[dashed] (A) -- (D);
  \draw[dashed] (P) -- (E);
  \draw[dashed] (B) -- (F);
  \draw[dashed] (F) -- (S);

  % 顶点标注
  \node[left] at (P) {$P$};
  \node[left] at (A) {$A$};
  \node[below left] at (B) {$B$};
  \node[below] at (C) {$C$};
  \node[right] at (D) {$D$};
  \node[below left] at (E) {$E$};
  \node[above right] at (F) {$F$};
  \node[right] at (S) {$S$};
\end{tikzpicture}
